% !TeX root = ../navier-stokes-manuscript.tex
% ============================================================
% 2.1 Vortex Stretching
% ============================================================

Axial stretching lengthens one narrow material vortex-tube segment while its material volume stays fixed, so the segment contracts across its axis, its rotating fluid moves inward, and its interior spins faster. The velocity field carrying this segment obeys the momentum equation
\[
  \underbrace{\partial_t u}_{\text{local acceleration}}
  +
  \underbrace{(u\cdot\nabla)u}_{\text{nonlinear transport}}
  =
  \underbrace{-\nabla p}_{\text{pressure}}
  +
  \underbrace{\nu\Delta u}_{\text{viscous diffusion}}.
\]
Taking its curl removes pressure and turns nonlinear transport into vorticity advection and stretching.

\subsection{Vortex Stretching}\label{sub:ThePerfectlyStretchingVortex}

The strain is locally uniform across the material segment and aligned with its axis. If \(e\) is that axis and \(L(t)\) is the segment's length, write
\[
  S:=\frac12\bigl(\nabla u+(\nabla u)^{\mathsf T}\bigr),
  \qquad
  Se=\sigma(t)e,
  \qquad
  \sigma(t)>0,
  \qquad
  \frac{\dot L(t)}{L(t)}
  =
  e\cdot Se
  =
  \sigma(t).
\]
Its material volume \(\mathcal V\) stays fixed. With effective transverse area \(A(t):=\mathcal V/L(t)\) and effective radius \(r_{\mathrm{eff}}(t):=\sqrt{A(t)/\pi}\),
\[
  \nabla\cdot u=0,
  \qquad
  \frac{d}{dt}(A L)=0,
  \qquad
  \frac{\dot A}{A}=-\sigma(t),
  \qquad
  \frac{\dot r_{\mathrm{eff}}}{r_{\mathrm{eff}}}
  =
  -\frac{\sigma(t)}{2}.
\]
The vorticity \(\omega:=\nabla\times u\) carried by this narrowing segment satisfies
\[
  \frac{D\omega}{Dt}
  =
  S\omega+\nu\Delta\omega,
  \qquad
  \frac{D}{Dt}
  :=
  \partial_t+u\cdot\nabla,
\]
and under the alignment \(\omega=|\omega|e\),
\[
  S\omega
  =
  \sigma(t)\omega,
  \qquad
  \frac12\frac{D|\omega|^2}{Dt}
  =
  \sigma(t)|\omega|^2
  +
  \nu\,\omega\cdot\Delta\omega.
\]
Aligned stretching contributes positively, while viscosity retains its sign; squared vorticity grows exactly where their full signed sum is positive. The energy identity proved in \Cref{prop:ClassicalKineticEnergyIdentity} and the antisymmetric part of \(\nabla u\) give
\[
  \norm{u(t)}_2
  \leq
  \norm{u_0}_2
  <\infty,
  \qquad
  \left|\frac12\bigl(\nabla u-(\nabla u)^{\mathsf T}\bigr)\right|_{\mathrm F}
  =
  \frac{|\omega|}{\sqrt2}
  \leq
  |\nabla u|_{\mathrm F}.
\]
Finite kinetic energy can coexist with concentration in the rotational, antisymmetric part of the velocity gradient. The exact comparison must consequently follow the full field to a finer width.

\divider{\NavierStokes{} Scaling}

At a time \(t_0<T_*\), let \(\ell:=r_{\mathrm{eff}}(t_0)\). For \(\lambda>1\), the full-field comparison at width \(\lambda^{-1}\ell\) is
\[
  u_\lambda(x,t)
  :=
  \lambda u(\lambda x,\lambda^2t),
  \qquad
  p_\lambda(x,t)
  :=
  \lambda^2p(\lambda x,\lambda^2t).
\]
On the corresponding rescaled interval, \(u_\lambda\) is another \NavierStokes{} solution. It is neither a complete solution nor a later state of the opening material segment. At the corresponding time \(\lambda^{-2}t_0\), its effective width is \(\lambda^{-1}\ell\), and
\[
\begin{aligned}
  \partial_tu_\lambda(x,t)
  &=
  \lambda^3(\partial_tu)(\lambda x,\lambda^2t),
  \\
  (u_\lambda\cdot\nabla)u_\lambda(x,t)
  &=
  \lambda^3\bigl((u\cdot\nabla)u\bigr)(\lambda x,\lambda^2t),
  \\
  \nabla p_\lambda(x,t)
  &=
  \lambda^3(\nabla p)(\lambda x,\lambda^2t),
  \\
  \nu\Delta u_\lambda(x,t)
  &=
  \lambda^3\nu(\Delta u)(\lambda x,\lambda^2t).
\end{aligned}
\]
Rescaling multiplies every momentum term by the same cubic factor, so nonlinear transport and viscosity remain matched.

\divider{Critical Height}

The momentum balance survives the change of width, while the Sobolev norms change with their order. Let \(\Lambda:=(-\Delta)^{1/2}\). Since
\[
  \widehat{u_\lambda}(\xi,t)
  =
  \lambda^{-2}\widehat u(\lambda^{-1}\xi,\lambda^2t),
\]
a change of variables gives
\[
  \norm{u_\lambda(t)}_{\dot H^s}^2
  =
  \lambda^{2s-1}
  \norm{u(\lambda^2t)}_{\dot H^s}^2.
\]
At \(s=0\), kinetic energy loses one power of \(\lambda\). At \(s=1\), first derivatives gain one power. At \(s=\tfrac12\), the power vanishes, leaving the scale-neutral level between them. Define
\[
  H(t)
  :=
  \frac12\norm{u(t)}_{\dot H^{1/2}}^2
  =
  \frac12\norm{\Lambda^{1/2}u(t)}_2^2.
\]
Writing \(H_\lambda\) for the same quantity evaluated on \(u_\lambda\),
\[
  \norm{u_\lambda(t)}_2^2
  =
  \lambda^{-1}\norm{u(\lambda^2t)}_2^2,
  \qquad
  H_\lambda(t)=H(\lambda^2t),
  \qquad
  \norm{\nabla u_\lambda(t)}_2^2
  =
  \lambda\norm{\nabla u(\lambda^2t)}_2^2.
\]
At corresponding rescaled times, kinetic energy decreases by one power, the first-derivative norm increases by one, and \(H\) stays unchanged by scale. This neutrality gives no direction to the time evolution of \(H\).

\divider{Nonlinear Production versus Viscous Dissipation}

At critical height, nonlinear transport contributes the signed production
\[
  \mathcal P_{1/2}[u](t)
  :=
  -\operatorname{Re}
  \pair
    {\Lambda^{1/2}u(t)}
    {\Lambda^{1/2}\bigl((u(t)\cdot\nabla)u(t)\bigr)}_{L^2}.
\]
Incompressibility removes pressure from the pairing, while viscosity contributes the negative dissipation term \(-\nu\norm{\Lambda^{3/2}u}_2^2\). Thus
\[
  H'(t)
  +
  \nu\norm{\Lambda^{3/2}u(t)}_2^2
  =
  \mathcal P_{1/2}[u](t).
\]
The height rises when nonlinear production exceeds viscous dissipation and falls when dissipation is larger. Under the same rescaling,
\[
  \mathcal P_{1/2}[u_\lambda](t)
  =
  \lambda^2\mathcal P_{1/2}[u](\lambda^2t),
  \qquad
  \nu\norm{\Lambda^{3/2}u_\lambda(t)}_2^2
  =
  \lambda^2\nu\norm{\Lambda^{3/2}u(\lambda^2t)}_2^2.
\]
Rescaling multiplies production and dissipation by the same quadratic factor. Neither gains ground, and the duration of a nonlinear transition remains undetermined.

\divider{The Heat Clock}

The heat equation \(v_t=\nu\Delta v\) supplies a linear viscous comparison. Its Fourier modes evolve by
\[
  \widehat v(\xi,t+\delta t)
  =
  e^{-\nu|\xi|^2\delta t}\widehat v(\xi,t).
\]
A scale-localized velocity structure of width \(r\) occupies Fourier frequencies with magnitude comparable to \(r^{-1}\). Its viscous change becomes order one when
\[
  \nu|\xi|^2\delta t\simeq1,
\]
so its linear heat comparison time is
\[
  \tau_\nu(r)
  :=
  \frac{r^2}{\nu}.
\]
Halving the width quarters this comparison time. More generally,
\[
  \tau_\nu(\lambda^{-1}r)
  =
  \lambda^{-2}\tau_\nu(r).
\]
This linear contraction does not determine the lifetime of a nonlinear structure.

\Needspace{18\baselineskip}
\divider{The Zeno Cascade}

The dyadic widths and their heat times are defined by
\[
  r_n:=2^{-n}r_0,
  \qquad
  \tau_n:=\frac{r_n^2}{\nu},
\]
respectively. The heat times shrink geometrically:
\[
  \tau_n
  =
  \frac{r_0^2}{\nu}4^{-n},
  \qquad
  \sum_{n=0}^{\infty}\tau_n
  =
  \frac{4r_0^2}{3\nu}
  <
  \infty.
\]
Their finite sum can fit inside a finite interval, but this arithmetic does not realize a \NavierStokes{} history. One velocity field would have to exhibit scale-localized structures at the widths
\[
  r_0>r_1>r_2>\cdots\longrightarrow0
\]
and at sampled times
\[
  t_0<t_1<t_2<\cdots\uparrow T_*<\infty,
  \qquad
  t_{n+1}-t_n\asymp\tau_n,
\]
with comparison constants independent of \(n\). Its remaining time would then satisfy
\[
  T_*-t_n
  \asymp
  \sum_{k=n}^{\infty}\tau_k
  =
  \frac{4r_n^2}{3\nu}.
\]
If one field realizes this history, that same field passes through successively smaller localized structures in gaps comparable to \(r_n^2/\nu\), and those gaps collapse as the sampled times approach the finite endpoint \(T_*\). This schedule gives no control of that field's first or successive Eulerian time derivatives at one fixed spatial point.
